\documentclass[a4paper, 11pt]{article}
%Enable URL highlighting
\usepackage{hyperref}
\usepackage{graphicx}
\usepackage[table,xcdraw]{xcolor}
\hypersetup{
    colorlinks=true,
    linkcolor=blue,
    urlcolor=blue,
}
\urlstyle{same}

%glossary
\usepackage[acronym,toc]{glossaries}
\makeglossaries
\setglossarystyle{altlistgroup}

\usepackage{xparse}
\DeclareDocumentCommand{\newdualentry}{ O{} O{} m m m m } {
  \newglossaryentry{gls-#3}{name={#5},text={#5\glsadd{#3}},
    description={#6},#1
  }
  \makeglossaries
  \newacronym[see={[Glossary:]{gls-#3}},#2]{#3}{#4}{#5\glsadd{gls-#3}}
}

\makeatletter
\renewcommand{\paragraph}{\@startsection{paragraph}{4}{0ex}%
    {-3.25ex plus -1ex minus -0.2ex}%
    {0.3ex plus 0.2ex}%
    {\normalfont\normalsize\bfseries}}
\makeatother

%page dimensions
\usepackage[a4paper,includeheadfoot, top=2.5cm, bottom=2.5cm, left=3.5cm, right=3cm]{geometry}

%Font settings
\usepackage{tgpagella}
\usepackage[T1]{fontenc}

%Support for German
\usepackage[utf8]{inputenc}
\usepackage[ngerman]{babel}

%number everything down to paragraphs
\setcounter{secnumdepth}{5}

%Enable Microtyping which improves justification
\usepackage{microtype}

%header
\setlength{\headheight}{15pt}
\usepackage{fancyhdr}
\pagestyle{fancy}
\lhead{\nouppercase{\leftmark}}
\rhead{\includegraphics[width=1cm]{pictures/ibzlogo.eps}}
\cfoot{\thepage}

%footnotes
\setlength{\skip\footins}{0.5cm}
\usepackage[bottom,hang]{footmisc}

%European style paragraphs
\setlength{\parskip}{\baselineskip}%
\setlength{\parindent}{0pt}%

%Make line spacing bigger
\linespread{1.3}

%Make whitespace in tables bigger
\renewcommand\arraystretch{1.5}

%Enable colours for tables
\usepackage[table]{xcolor}
\usepackage{longtable}

%Needed to rotate graphics and tables
\usepackage{rotating}
\usepackage{pdflscape}

%Support for images
\usepackage{graphicx}
\usepackage{float}

%Material theme colours
\definecolor{red}{HTML}{F44336}
\definecolor{pink}{HTML}{E91E63}
\definecolor{purple}{HTML}{9C27B0}
%\definecolor{blue}{HTML}{2196F3}
\definecolor{brown}{HTML}{795548}
\definecolor{cyan}{HTML}{00BCD4}
\definecolor{darkgray}{HTML}{616161}
\definecolor{gray}{HTML}{9E9E9E}
\definecolor{lightgray}{HTML}{E0E0E0}
\definecolor{lime}{HTML}{CDDC39}
\definecolor{olive}{HTML}{827717}
\definecolor{orange}{HTML}{FF9800}
\definecolor{teal}{HTML}{009688}
\definecolor{yellow}{HTML}{FFEB3B}
\definecolor{green}{HTML}{388E3C}

%Support for Code Snippets and Syntax Highlighting
\usepackage{listings}
\usepackage{color}
\usepackage{sourcecodepro}
\lstset{
  basicstyle=\ttfamily,
  language=bash,
  breaklines,
  literate={ö}{{\"o}}1
           {ä}{{\"a}}1
           {ü}{{\"u}}1,
  captionpos=b,                    % sets the caption-position to bottom
  commentstyle=\color{green},    % comment style
  escapeinside={\%*}{*)},          % if you want to add LaTeX within your code
  keywordstyle=\color{blue},       % keyword style
  stringstyle=\color{purple},     % string literal style
  showstringspaces=false,          % Removes the strange symboles where spaces are
}

% Listing with a box around it
\usepackage[most]{tcolorbox}
\newtcblisting{sexylisting}[2][]{
  sharp corners,
  fonttitle=\bfseries,
  colframe=gray,
  listing only,
  listing options={basicstyle=\ttfamily,language=bash},
  title=\thetcbcounter #2, #1
}

\usepackage[citestyle=verbose,bibstyle=numeric,sorting=none,backend=bibtex8]{biblatex}


\bibliography{bib,andreas_general,andreas}
%\include{glossary}
\begin{document}

\begin{titlepage}
    \centering
    %\includegraphics[width=0.15\textwidth]{example-image-1x1}\par\vspace{1cm}
    {\scshape\LARGE IBZ-Schulen AG\par}
    %\vspace{0.5cm}
    {\scshape\Large Aarau\par}
    \vspace{1.5cm}
    {\huge\bfseries Casestudy Webtechnologien\par}
    \vspace{2cm}
    {\Large\itshape Hörler Ivan, Zweili Andreas, TI-AR-17W\par}
    \vfill

% Bottom of the page
    {\large \today\par}
\thispagestyle{empty}
\end{titlepage}


\newpage
\renewcommand{\abstractname}{Management Summary}
\begin{abstract}
Dies ist die Dokumentation für die zweite Case Study im Fach
Webtechnologien von Ivan Hörler und Andreas Zweili, welche diese im
Rahmen ihres 5. Semesters an der IBZ Schule in Aarau erarbeiteten. Die
Case Study behandelt dabei das Erstellen eines Web-Shops und der dafür
gewählten Werkzeuge, die Projektplanung sowie die dabei aufgetretenen
Probleme.

Zusätzlich sollen auch die Erfahrungen der Studenten im Zusammenhang
des verwendeten Frameworks Django aufgezeigt werden. Dieses ist nicht
Teil des Kurikulums, weshalb diese Arbeit interessante, zusätzliche
Möglichkeiten im Bereich der Entwicklung von Webapplikationen
aufzeigen kann.
\end{abstract}

\newpage
\microtypesetup{protrusion=false} % disables protrusion locally in the document
\tableofcontents % prints Table of Contents
\microtypesetup{protrusion=true} % enables protrusion
\newpage


\include{doku}

\newpage
\nocite{*}
\printbibliography[heading=bibintoc]

\newpage
\microtypesetup{protrusion=false}
\listoffigures
\microtypesetup{protrusion=true}
\
\newpage

\microtypesetup{protrusion=false}
\listoftables
\microtypesetup{protrusion=true}
\
\newpage


%\printglossaries
\end{document}


%%% Local Variables:
%%% mode: latex
%%% TeX-master: t
%%% End:
